\clearpage
\section{Chapter 3: Cosmic Evolution}

\subsection{Questions}

\begin{enumerate}[label=3.\arabic*]
  \item \label{q:3.1} \textbf{Continuity Equation with EoS}

  Starting from
  \[
    \dot{\rho}+3H(\rho+p)=0,
  \]
  and using $p=w\rho$ with constant $w$:
  \begin{enumerate}[label=(\alph*)]
    \item Derive $\rho(a)\propto a^{-3(1+w)}$.
    \item Give the scaling for dust, radiation, and cosmological constant.
    \item \textit{(Toy assumption for scaling practice)} If $\rho_m=\rho_r$ today, what is $\rho_m/\rho_r$ at $z=9$?
  \end{enumerate}

  \item \label{q:3.2} \textbf{Flat $\Lambda$CDM Expansion Rate}

  Assume a flat universe ($k=0$), negligible radiation today, and
  \[
    \Omega_{m0}=0.30,\qquad \Omega_{\Lambda0}=0.70.
  \]
  \begin{enumerate}[label=(\alph*)]
    \item Write $H(z)$ in terms of $H_0$.
    \item Compute $H(z=1)/H_0$.
    \item \textit{(Hint: from acceleration equation)} show that today
    \[
      q_0 \equiv -\frac{\ddot a a}{\dot a^2}\Big|_0 = \frac{1}{2}\Omega_{m0}-\Omega_{\Lambda0},
    \]
    and evaluate $q_0$.
  \end{enumerate}

  \item \label{q:3.3} \textbf{Newtonian Limit with Cosmological Constant}

  Use
  \[
    F(r) = -\frac{GMm}{r^2}+\frac{1}{3}\Lambda r m.
  \]
  \begin{enumerate}[label=(\alph*)]
    \item Find the radius $R_\mathrm{eq}$ where attraction and repulsion are equal.
    \item Estimate $R_\mathrm{eq}$ for the Sun using $M_\odot=1.989\times10^{30}\,\mathrm{kg}$,
    $G=6.674\times10^{-11}$ SI,
    $\Lambda=1.1\times10^{-52}\,\mathrm{m^{-2}}$.
    \item Estimate $F_\Lambda/F_G$ at $r=1\,\mathrm{AU}$.
  \end{enumerate}

  \item \label{q:3.4} \textbf{Redshift-Form $\Lambda$CDM and Equality Epochs}

  Assume a near-flat universe with
  \[
    \Omega_{m0}=0.315,\qquad
    \Omega_{r0}=9.4\times10^{-5},\qquad
    \Omega_{\Lambda0}=0.685,
  \]
  and ignore curvature in this exercise.
  \begin{enumerate}[label=(\alph*)]
    \item Write $E(z)\equiv H(z)/H_0$.
    \item Find the redshift $z_{\mathrm{eq}}^{(mr)}$ where matter and radiation are equal.
    \item Find the redshift $z_{\mathrm{eq}}^{(m\Lambda)}$ where matter and cosmological constant are equal.
  \end{enumerate}

  \item \label{q:3.5} \textbf{Onset of Accelerated Expansion}

  Assume a flat universe with negligible radiation and
  \[
    \Omega_{m0}=0.315,\qquad \Omega_{\Lambda0}=0.685.
  \]
  \begin{enumerate}[label=(\alph*)]
    \item Starting from
    \[
      q(z)=\frac{\frac12\Omega_{m0}(1+z)^3-\Omega_{\Lambda0}}{\Omega_{m0}(1+z)^3+\Omega_{\Lambda0}},
    \]
    find the transition redshift $z_{\mathrm{acc}}$ where expansion changes from deceleration to acceleration.
    \item Compare $z_{\mathrm{acc}}$ with the matter--$\Lambda$ equality redshift
    \[
      z_{\mathrm{eq}}^{(m\Lambda)}=\left(\frac{\Omega_{\Lambda0}}{\Omega_{m0}}\right)^{1/3}-1,
    \]
    and explain physically why the two redshifts are different.
  \end{enumerate}
\end{enumerate}

\bigskip
\bigskip
\bigskip

\subsection{Solutions}

\subsubsection*{Solution 3.1: Continuity Equation with EoS}

\begin{enumerate}[label=(\alph*)]
  \item
  \[
    \dot\rho + 3\frac{\dot a}{a}(\rho+p)=0,
    \qquad p=w\rho
  \]
  gives
  \[
    \dot\rho + 3(1+w)\frac{\dot a}{a}\rho=0
    \Rightarrow \frac{d\rho}{\rho}=-3(1+w)\frac{da}{a}.
  \]
  Integrating:
  \[
    \rho\propto a^{-3(1+w)}.
  \]

  \item
  \[
    \rho_m\propto a^{-3},\qquad
    \rho_r\propto a^{-4},\qquad
    \rho_\Lambda=\text{constant}.
  \]

  \item
  \[
    \frac{\rho_m}{\rho_r}\propto a,
    \quad a=\frac{1}{1+z}.
  \]
  At $z=9$, $a=0.1$, so
  \[
    \left(\frac{\rho_m}{\rho_r}\right)_{z=9}=0.1\left(\frac{\rho_m}{\rho_r}\right)_0=0.1.
  \]
  This result follows from the toy assumption in the question; it is mathematically correct even though the assumption is not physically realistic for the actual Universe today.
\end{enumerate}

\subsubsection*{Solution 3.2: Flat $\Lambda$CDM Expansion Rate}

\begin{enumerate}[label=(\alph*)]
  \item For flat universe and negligible radiation:
  \[
    \frac{H^2(z)}{H_0^2}=\Omega_{m0}(1+z)^3+\Omega_{\Lambda0}.
  \]

  \item At $z=1$:
  \[
    \frac{H^2}{H_0^2}=0.30\times 2^3 + 0.70 = 3.10
    \Rightarrow \frac{H}{H_0}=\sqrt{3.10}\approx 1.76.
  \]

  \item From
  \[
    \frac{\ddot a}{a}=-\frac{4\pi G}{3}(\rho+3p)+\frac{\Lambda}{3},
  \]
  and using $p_m\approx 0$, $p_\Lambda=-\rho_\Lambda$,
  \[
    q_0 = \frac{1}{2}\Omega_{m0}-\Omega_{\Lambda0}
    = 0.15-0.70 = -0.55.
  \]
  Negative $q_0$ means accelerated expansion. (Here radiation is neglected, consistent with the question statement.)
\end{enumerate}

\subsubsection*{Solution 3.3: Newtonian Limit with Cosmological Constant}

\begin{enumerate}[label=(\alph*)]
  \item Set magnitudes equal:
  \[
    \frac{GMm}{R_\mathrm{eq}^2}=\frac{1}{3}\Lambda R_\mathrm{eq}m
    \Rightarrow
    R_\mathrm{eq}=\left(\frac{3GM}{\Lambda}\right)^{1/3}.
  \]

  \item
  \[
    R_\mathrm{eq}\approx \left(\frac{3(6.674\times10^{-11})(1.989\times10^{30})}{1.1\times10^{-52}}\right)^{1/3}
    \approx 1.5\times10^{24}\,\mathrm{m}.
  \]
  In parsec ($1\,\mathrm{pc}=3.086\times10^{16}\,\mathrm{m}$):
  \[
    R_\mathrm{eq}\approx 5\times10^7\,\mathrm{pc} \;(\sim 50\,\mathrm{Mpc}),
  \]
  far beyond Solar System scales.

  \item
  \[
    \frac{F_\Lambda}{F_G}=
    \frac{\frac13\Lambda r m}{\frac{GMm}{r^2}}=
    \frac{\Lambda r^3}{3GM}.
  \]
  At $r=1\,\mathrm{AU}=1.496\times10^{11}\,\mathrm{m}$,
  \[
    \frac{F_\Lambda}{F_G}\approx 3\times10^{-23}.
  \]
  So $\Lambda$ is negligible for planetary dynamics.
\end{enumerate}

\subsubsection*{Solution 3.4: Redshift-Form $\Lambda$CDM and Equality Epochs}

\begin{enumerate}[label=(\alph*)]
  \item Ignoring curvature,
  \[
    E(z)=\sqrt{\Omega_{m0}(1+z)^3+\Omega_{r0}(1+z)^4+\Omega_{\Lambda0}}.
  \]
  With given numbers:
  \[
    E(z)=\sqrt{0.315(1+z)^3+9.4\times10^{-5}(1+z)^4+0.685}.
  \]

  \item Matter--radiation equality:
  \[
    \Omega_{m0}(1+z)^3=\Omega_{r0}(1+z)^4
    \Rightarrow
    1+z_{\mathrm{eq}}^{(mr)}=\frac{\Omega_{m0}}{\Omega_{r0}}.
  \]
  Therefore
  \[
    1+z_{\mathrm{eq}}^{(mr)}=\frac{0.315}{9.4\times10^{-5}}\approx 3351
    \Rightarrow
    z_{\mathrm{eq}}^{(mr)}\approx 3.35\times10^3.
  \]

  \item Matter--$\Lambda$ equality:
  \[
    \Omega_{m0}(1+z)^3=\Omega_{\Lambda0}
    \Rightarrow
    1+z_{\mathrm{eq}}^{(m\Lambda)}=\left(\frac{\Omega_{\Lambda0}}{\Omega_{m0}}\right)^{1/3}.
  \]
  Hence
  \[
    1+z_{\mathrm{eq}}^{(m\Lambda)}=\left(\frac{0.685}{0.315}\right)^{1/3}\approx 1.295
    \Rightarrow
    z_{\mathrm{eq}}^{(m\Lambda)}\approx 0.295.
  \]
\end{enumerate}

\subsubsection*{Solution 3.5: Onset of Accelerated Expansion}

\begin{enumerate}[label=(\alph*)]
  \item The transition occurs at $q(z_{\mathrm{acc}})=0$, so the numerator must vanish:
  \[
    \frac12\Omega_{m0}(1+z_{\mathrm{acc}})^3-\Omega_{\Lambda0}=0.
  \]
  Hence
  \[
    (1+z_{\mathrm{acc}})^3 = \frac{2\Omega_{\Lambda0}}{\Omega_{m0}}
    \Rightarrow
    z_{\mathrm{acc}} = \left(\frac{2\Omega_{\Lambda0}}{\Omega_{m0}}\right)^{1/3}-1.
  \]
  Numerically,
  \[
    z_{\mathrm{acc}}=\left(\frac{2\times0.685}{0.315}\right)^{1/3}-1
    = (4.349)^{1/3}-1
    \approx 1.633-1
    \approx 0.633.
  \]

  \item Matter--$\Lambda$ equality is
  \[
    z_{\mathrm{eq}}^{(m\Lambda)}=\left(\frac{\Omega_{\Lambda0}}{\Omega_{m0}}\right)^{1/3}-1
    \approx 0.295.
  \]
  Therefore
  \[
    z_{\mathrm{acc}}\approx 0.633 > z_{\mathrm{eq}}^{(m\Lambda)}\approx 0.295.
  \]
  This is physically expected: acceleration requires
  \[
    \rho+3p<0.
  \]
  For matter plus $\Lambda$,
  \[
    \rho+3p = \rho_m + \rho_\Lambda + 3(0-\rho_\Lambda)=\rho_m-2\rho_\Lambda.
  \]
  So the threshold is not $\rho_m=\rho_\Lambda$, but
  \[
    \rho_m=2\rho_\Lambda,
  \]
  which happens at higher redshift.
\end{enumerate}


% ─────────────────────────────────────────────────────────────────────────────
